% Vietnamese translation for OI Public Library
% Source-File: IOI中国国家候选队论文/2008/Day1/9.周小博《浅谈信息学竞赛中的区间问题》/浅谈信息学竞赛中的区间问题.ppt
\documentclass[11pt]{article}
\usepackage{oipl-vn}
\usepackage{tikz}
\usetikzlibrary{arrows.meta,calc}

\newcommand{\Oh}{\mathcal{O}}
\newcommand{\depth}{d}

\title{Bản trình chiếu: Các bài toán đoạn trong thi tin học}
\author{Zhou Xiaobo}
\date{2008}

\begin{document}
\maketitle

\origtitle{Bản trình chiếu: Bàn về các bài toán đoạn trong thi tin học}
\sourcepath{IOI China 2008 / Day 1 / Zhou Xiaobo / interval problems.ppt}

\section{Mục tiêu của slide}

PPT của Zhou Xiaobo là một bản tổng kết theo mô hình. Thay vì bắt đầu từ tên
cấu trúc dữ liệu, slide bắt đầu từ câu hỏi: đối tượng nào có thể xem là một
đoạn trên trục, và điều kiện của đề tương ứng với thao tác gì trên các đoạn?
Sáu nhóm chính trong slide khớp với bài viết:
\begin{enumerate}
  \item chọn nhiều đoạn đôi một không giao nhau nhất;
  \item gán đoạn vào nhiều tài nguyên;
  \item đặt đoạn có hạn chót;
  \item phủ một đoạn mục tiêu bằng ít đoạn nhất;
  \item các biến thể có trọng số;
  \item quan hệ giữa đoạn và điểm.
\end{enumerate}

\section{Chọn nhiều đoạn nhất}

Mô hình cơ sở: cho \(n\) đoạn, chọn nhiều đoạn đôi một không giao nhau nhất.
Thuật toán tham lam sắp theo đầu phải tăng dần, rồi chọn đoạn đầu tiên không
giao với đoạn đã chọn cuối cùng.

\begin{center}
\begin{tikzpicture}[x=0.95cm,y=0.55cm, every node/.style={font=\small}]
  \draw[-{Latex[length=2mm]}] (0,0) -- (7.2,0);
  \foreach \a/\b/\y/\c in {0.4/2.0/1/blue,1.1/3.0/2/gray,2.2/3.4/1/blue,3.1/5.4/2/gray,3.7/4.7/1/blue,5.0/6.6/1/blue}
    \draw[line width=2pt,\c!70] (\a,\y) -- (\b,\y);
  \node[right] at (7.2,0) {thời gian};
\end{tikzpicture}
\end{center}

Chứng minh trong slide là chứng minh đổi chỗ chuẩn. Nếu đoạn đầu tiên của một
nghiệm tối ưu kết thúc sau đoạn tham lam chọn, ta thay nó bằng đoạn tham lam;
số đoạn không giảm và phần còn lại còn nhiều không gian hơn. Lặp lại cho từng
vị trí, thu được nghiệm tối ưu bắt đầu giống thuật toán. Độ phức tạp
\(\Oh(n\log n)\) do bước sắp xếp.

\section{Nhiều tài nguyên và độ sâu}

Mô hình thứ hai: mọi đoạn phải được gán vào tài nguyên, trên cùng một tài
nguyên các đoạn không giao nhau, và cần dùng ít tài nguyên nhất. Độ sâu
\(\depth\) của tập đoạn là số đoạn lớn nhất cùng phủ một điểm.

Rõ ràng cần ít nhất \(\depth\) tài nguyên. Ngược lại, quét các đoạn theo đầu
trái; khi xử lý một đoạn \(I\), mọi đoạn trước đó giao với \(I\) đều chứa đầu
trái của \(I\). Vì tại điểm ấy có nhiều nhất \(\depth\) đoạn chồng nhau, trong
\(\depth\) tài nguyên luôn còn một tài nguyên trống cho \(I\). Do đó số tài
nguyên tối thiểu đúng bằng độ sâu.

Cài đặt thường dùng hàng đợi ưu tiên hoặc cây cân bằng lưu thời điểm kết thúc
cuối cùng của mỗi tài nguyên. Khi cần nhận một đoạn mới, chọn tài nguyên có
đầu phải phù hợp nhất; nếu không có thì mở tài nguyên mới. Tư tưởng này cũng
xuất hiện trong ví dụ phân bến cảng trong bài viết.

\section{Đoạn có hạn chót}

Mô hình thứ ba không cố định vị trí đoạn. Mỗi công việc có độ dài \(p_i\) và
hạn chót \(d_i\); cần đặt các đoạn liên tiếp sao cho độ trễ lớn nhất
\[
  L_{\max}=\max_i(r_i-d_i)
\]
nhỏ nhất. Slide đưa ra quy tắc kinh điển: sắp theo hạn chót tăng dần, rồi đặt
liên tiếp không để khoảng trống.

Hai quan sát đủ cho chứng minh. Thứ nhất, trong một nghiệm tối ưu có thể dời
mọi phần sau một khoảng trống sang trái, nên tồn tại nghiệm tối ưu không có
khoảng trống. Thứ hai, nếu hai đoạn kề nhau bị nghịch thế theo hạn chót, hoán
đổi chúng không làm tăng độ trễ cực đại. Lặp hoán đổi sẽ cho thứ tự theo hạn
chót.

Ví dụ trong bài viết như \textit{Pendant} và \textit{Loan Scheduling} cho thấy
mô hình hạn chót thường cần thêm DP, cây đoạn hoặc DSU khi có chọn tập con,
trọng số, hoặc nhiều tài nguyên.

\section{Phủ đoạn tối thiểu}

Mô hình thứ tư: chọn ít đoạn nhất để phủ kín \([s,t]\). Thuật toán tham lam
giữ điểm đã phủ tới \(x\). Trong tất cả các đoạn có đầu trái không vượt quá
\(x\), chọn đoạn có đầu phải xa nhất, rồi cập nhật \(x\).

Sau \(k\) bước, điểm xa nhất mà thuật toán phủ được không nhỏ hơn điểm xa nhất
mà bất kỳ cách chọn \(k\) đoạn nào có thể bảo đảm theo cùng quy tắc xuất phát.
Vì vậy khi thuật toán đã phủ tới \(t\), không có nghiệm dùng ít đoạn hơn. Nếu
tại một bước không có đoạn nào bắt đầu trước hoặc tại \(x\), đoạn mục tiêu
không thể phủ kín.

\section{Có trọng số}

Khi mỗi đoạn có trọng số hoặc lợi nhuận, các quy tắc tham lam cơ sở thường
không còn đủ. Slide nhóm các biến thể thành hai hướng:
\begin{itemize}
  \item chọn các đoạn không giao nhau với tổng trọng số lớn nhất;
  \item phủ hoặc gán đoạn sao cho chi phí nhỏ nhất hay lợi nhuận lớn nhất.
\end{itemize}
Với chọn đoạn có trọng số, cách chuẩn là sắp theo đầu phải, tính
\(p(i)\) là đoạn cuối cùng kết thúc trước khi đoạn \(i\) bắt đầu, rồi dùng
\[
  f_i=\max(f_{i-1}, f_{p(i)}+w_i).
\]
Tìm \(p(i)\) bằng tìm kiếm nhị phân hoặc cây cân bằng, tổng thời gian
\(\Oh(n\log n)\).

Với phủ đoạn có trọng số, trạng thái thường là điểm đã phủ xa nhất hoặc đoạn
cuối cùng được chọn; cần nén tọa độ nếu đầu mút lớn. Điểm slide muốn nhấn
mạnh là: trọng số không phá mô hình đoạn, nhưng thường chuyển lời giải từ
tham lam sang quy hoạch động có tối ưu hóa.

\section{Quan hệ giữa đoạn và điểm}

Nhóm cuối cùng không chỉ chọn đoạn, mà còn hỏi quan hệ giữa đoạn và điểm: một
điểm bị bao bởi bao nhiêu đoạn, chọn điểm để đâm tất cả các đoạn, hoặc chọn
đoạn để bao nhiều điểm nhất. Với dữ liệu tĩnh, quét sự kiện là công cụ đầu
tiên. Với truy vấn động, ta dùng cây Fenwick, cây đoạn, cây cân bằng hoặc nén
tọa độ tùy thao tác.

Ví dụ slide nhắc tới các công thức dạng \(f[i]\), các đoạn \([M,E]\),
\([M,T2_i]\), và thao tác \(\Oh(\log N)\). Đây là dấu hiệu của cách làm chung:
sắp xếp đầu mút, biến quan hệ giao/phủ thành truy vấn tiền tố hoặc truy vấn
đoạn, rồi cập nhật DP bằng cấu trúc dữ liệu.

\section{Cách nhận diện mô hình đoạn}

Khi đề bài nói về thời gian bắt đầu--kết thúc, hạn chót, vùng phủ trên một
trục, tài nguyên không được dùng đồng thời, hoặc chọn tập hoạt động tương
thích, nên thử đưa đối tượng về đoạn. Sau đó hỏi tiếp:
\begin{itemize}
  \item ta chọn nhiều nhất, ít nhất, hay tối ưu trọng số?
  \item đoạn có vị trí cố định hay có thể dời?
  \item các đoạn được phép chạm đầu mút không?
  \item có truy vấn động hay tất cả dữ liệu đã biết trước?
\end{itemize}
Những câu hỏi này quyết định giữa tham lam, DP, quét sự kiện và cấu trúc dữ
liệu. Đó cũng là thông điệp thực dụng của bản trình chiếu: mô hình đoạn rất
quen, nhưng từng biến thể cần chứng minh và cài đặt riêng.

\end{document}
